\section{Resumen}
La transcripción automática de música (Automatic Music Transcription, AMT) es un área de investigación dentro del campo de Music Information Retrieval (MIR) cuyo objetivo consiste en transformar señales de audio musical en representaciones simbólicas como notas, acordes o partituras. Este problema resulta especialmente complejo en instrumentos polifónicos como la guitarra acústica, donde múltiples notas se producen simultáneamente generando espectros armónicos superpuestos.

En este trabajo se presenta un enfoque para la detección e identificación automática de acordes de guitarra acústica a partir de grabaciones de audio. La metodología propuesta se basa en el procesamiento espectral de la señal y en el uso de modelos de aprendizaje automático para reconocer patrones armónicos asociados a distintos acordes musicales.

Este estudio busca contribuir al desarrollo de herramientas computacionales que faciliten el análisis automático de interpretaciones musicales y la generación de representaciones simbólicas a partir de señales de audio~\cite{benetos2019automatic}.